\label{chap:83-14}

\section{Derivation of the channels \(90/120\) through inclusion of hexad--octad incidences}

With the dodecad and the alignment of the
\cref{p12-83-c17-s1:sec:w12-w24-elevacion} fixed, let \(H\) be a hexad in the HMT coordinates,
let \(\ell(H)\subset D\) be its image, and let
\(O_H=\iota_D(\ell(H))\) be its unique lift. Marking a point
and a pair of points of the residual hexad yields
\[
F_6(H)=\ell(H)\times\binom{\ell(H)}2,
\qquad
|F_6|=6\binom62=90.
\]
If the marked point can traverse the octad:
\[
F_8(H)=O_H\times\binom{\ell(H)}2,
\qquad
|F_8|=8\binom62=120.
\]
The inclusion \(\ell(H)\hookrightarrow O_H\) induces the incidence morphism
\[
 \mathcal I_{6\to8}:F_6(H)\hookrightarrow F_8(H).
\]
Normalizing the difference by the chosen
factor \(24\binom62\) yields
\[
\boxed{
\frac{90-120}{24\binom62}
=\frac{6-8}{24}
=-\frac1{12}.}
\]
The equality is a normalized defect of the combinatorial interface \(6\to8\). The factor \(24\binom62\) forms part of the definition of this normalization; the inclusion alone determines the difference \(-30\). The operator \(\mathcal I_{6\to8}\) has marked incidence configurations as its domain; it is neither the transverse extractor of the dodecaphase state nor the family of angular moments.

\section{Transport of \(-1/12\) to the channels \(90/120\) in \(5\) algebraic realizations}

In addition to the normalized defect of the hexad--octad inclusion,
five operator realizations of the same rational number appear. Their domains are
specified separately; equality of values does not identify the
operators without an additional intertwiner.

\subsection{Difference of dodecaphasic periodicities}

One tooth of \(\Cdoce\) measures \(30^\circ\). Then
\[
\frac{30}{120}-\frac{30}{90}
=\frac14-\frac13
=-\frac1{12}.
\]
It is the oriented difference between the periods of \(90^\circ\) and
\(120^\circ\) in this realization.

\subsection{Pseudoinverse of the Gram Operator on the Live Component}

In the exact jump \(60\to81\), two fibers have sizes \(0\) and
\(12\). The Gram operator is
\[
K_{60,81}=\operatorname{diag}(0,12).
\]
On the nonzero component of its range, the pseudoinverse is \(1/12\); with the negative
boundary orientation,
\[
E_{\rm surv}=-K_{60,81}^{+}
\]
yields \(-1/12\).

\subsection{2nd difference between triadic scales}

Define
\[
T_L=\sum_{n=1}^{L-1}n(L-n)=\frac{L(L^2-1)}6,
\qquad
a(L)=-\frac{T_L}{L^2}=-\frac L6+\frac1{6L}.
\]
The first scale subtraction is
\[
A_1(L)=a(L)-\frac{a(3L)}3=\frac4{27L},
\]
and the second
\[
C(L)=LA_1(L)-\frac{3L}{3}A_1(3L)=\frac8{81}.
\]
The residue \(8/81\) is independent of \(L\). To obtain
\(-1/12\), the declared normalization is applied
\[
R(L)=-\frac{27}{32}C(L)=-\frac1{12}.
\]
The calculation derives \(8/81\); it does not by itself derive the factor \(-27/32\).

\subsection{Composition of modular transformations and preservation of the signed state}

The eta function satisfies
\[
\frac{\eta(\tau)}{\eta(3\tau)}
=q^{-1/12}\prod_{n\ge1}(1+q^n+q^{2n})^{-1}.
\]
The exponent comes from \(1/24-3/24=-1/12\). Taking twelve copies:
\[
t_3(q)=\left(\frac{\eta(\tau)}{\eta(3\tau)}\right)^{12}
=q^{-1}-12+54q-76q^2-243q^3+\cdots.
\]
The coefficient \(54\) coincides with the APP defect\@. The rational identity
\[
j=\frac{(t_3+27)(t_3+243)^3}{t_3^3}
\]
establishes the connection with the classical modular branch.

\subsection{Defect of cyclic difference projectors}

Let \(S\) be the cyclic shift on \(\QQ^{12}\) and define
\[
 D_r=I-S^r.
\]
The kernel of \(D_r\) consists of the vectors constant on each orbit
of \(S^r\); therefore,
\[
 \dim\ker D_r=\gcd(12,r).
\]
Let \(\Pi_{\ker D_3}\) and \(\Pi_{\ker D_4}\) be the rational projectors onto
\(\ker D_3\) and \(\ker D_4\), and let
\(\tau(T)=\operatorname{Tr}(T)/12\) be the normalized trace. Then
\[
 \boxed{
 \tau(\Pi_{\ker D_3}-\Pi_{\ker D_4})
 =\frac{\operatorname{rank}\Pi_{\ker D_3}
       -\operatorname{rank}\Pi_{\ker D_4}}{12}
 =\frac{3-4}{12}
 =-\frac1{12}.}
\]
This realization expresses the rational as the transverse defect of steps
three and four in the twelve-phase cycle.

\begin{remark}[Domains of the realizations]
The APP defect \(54\), the eta exponent \(-1/12\), the combinatorial
defect of the hexad--octad lift, and the
modular identities once \(t_3\) has been fixed are exact. The value
of the pseudoinverse, the normalized extractor, and the projector defect are likewise exact.
Each equality preserves its domain and its corresponding map.
\end{remark}


\section{Barbero--Immirzi parameter: area operator, spectrum,
channel intertwiner, and modular entropy}
\label{p12-83-c14-s2:sec:area-entrelazamiento-barbero-rev11}
\index{Barbero--Immirzi!area operator}
\index{modular entropy!modular Hamiltonian}

The HMT identification of the functional \(\Gamma_{6\to8}\) with the
Barbero--Immirzi coordinate is carried out on the same triadic boundary that supports the
incidence (90/120). The bulk and its boundary sheet have the same cardinality,
\[
 \mathcal B=(\mathbb Z/9\mathbb Z)^3,
 \quad |\mathcal B|=729,
 \qquad
 \partial\mathcal B=(\mathbb Z/27\mathbb Z)^2,
 \quad |\partial\mathcal B|=729,
\]
by regrouping six trits into three nonadic coordinates or two
coordinates of order twenty-seven.

\begin{theorem}[Areal Law of Triadic Cutoff]
\label{p12-83-c14-s2:thm:ley-areal-corte-triadico-rev11}
In the cubic graph \(N\times N\times N\) with unit capacity, the minimum
cut between two opposite faces has capacity \(N^2\). For \(N=9\),
\[
 \mathcal A_{\min}=81,
 \qquad S_{\rm tri}=\mathcal A_{\min}\log3=81\log3.
\]
\end{theorem}
\begin{proof}
The \(N^2\) lines parallel to the normal direction are edge-disjoint
paths, so every cut intercepts them. A transverse layer contains
exactly \(N^2\) edges and attains the bound. Each triadic link contributes
\(\log3\) to the entropy of a maximally mixed state.
\end{proof}

On the boundary torus \(27\times27\), let \(A\) be the canonical block
\(9\times9\). Its boundary decomposes exactly into
\begin{equation}
 \partial A=\partial A_{90}\sqcup\partial A_{120},
 \qquad |\partial A_{90}|=|\partial A_{120}|=18.
 \label{p12-83-c14-s2:eq:corte-90-120-rev11}
\end{equation}
The horizontal links represent channel (90), and the vertical ones
channel (120). This decomposition turns the pair of incidences into an
operator decomposition of the boundary.

The areal entropy \(S_{\rm tri}=81\log3\) belongs to the maximally
mixed state of the cut links. The bilayer entropy
\(S_{\rm bi}(y)\) of the \cref{def:entropia-bicapa-rev6} belongs to the
normalized distribution between the two sheets; the two objects retain
distinct domains.

Let \(N=\operatorname{diag}(0,1,2)\) be the number operator of a tritic link.
The areal realization declares the positive scale
\[
 \theta_{\rm clk}=\frac{C^*}{6}\frac{\pi}{180},
 \qquad
 a_0=\frac{1+2\cos(10^\circ)}{3}\,\theta_{\rm clk},
 \qquad \gamma_{\rm HMT}=\Gamma_{6\to8}.
\]
The formula for \(a_0\) constitutes the starting structure of this
realization; once that scale is fixed, the operator and its spectrum follow
exactly.

\begin{definition}[Area Operator HMT]
On the boundary sheet
\(\mathcal H_{\partial}=\bigotimes_{e\in\partial A}\mathbb C^3\), define
\begin{equation}
 \widehat{\mathcal A}_{\rm HMT}
 =\gamma_{\rm HMT}a_0\sum_{e\in\partial A}N_e.
 \label{p12-83-c14-s2:eq:operador-area-hmt-rev11}
\end{equation}
Its spectrum in the canonical cut is
\[
 \operatorname{Spec}\widehat{\mathcal A}_{\rm HMT}
 =\{n\gamma_{\rm HMT}a_0:0\leq n\leq72\},
\]
with multiplicities given by the coefficients of
\((1+z+z^2)^{36}\). Prolongation through compatible cuts produces the
protected family \(\mathcal A_n^{\rm prot}=n\gamma_{\rm HMT}a_0\).
\end{definition}

The intertwiner that fixes normalization is constructed in the subspace of
collective channels. Let \(u_{90},u_{120}\) be the uniform unit vectors
of the hexadic and octadic incidences, and let \(v_{90},v_{120}\) be the
uniform unit vectors of the two sets of links of
\eqref{p12-83-c14-s2:eq:corte-90-120-rev11}. The map
\begin{equation}
 \begin{aligned}
 \mathcal J_{6\to8}:{}
 &\operatorname{span}\{u_{90},u_{120}\}
 \longrightarrow
 \operatorname{span}\{v_{90},v_{120}\},\\
 &\mathcal J_{6\to8}u_{90}=v_{90},
 \qquad
 \mathcal J_{6\to8}u_{120}=v_{120}
 \end{aligned}
 \label{p12-83-c14-s2:eq:entrelazador-area-barbero-rev11}
\end{equation}
is an isometry between these collective subspaces. If
\[
 D_{\rm inc}=90|u_{90}\rangle\langle u_{90}|
 +120|u_{120}\rangle\langle u_{120}|,
 \qquad
 D_{\rm area}=90|v_{90}\rangle\langle v_{90}|
 +120|v_{120}\rangle\langle v_{120}|,
\]
then
\[
 \boxed{\mathcal J_{6\to8}D_{\rm inc}
 =D_{\rm area}\mathcal J_{6\to8}.}
\]
In particular, the primitive combination \(12:-1\) acts
with the same coefficients before and after \(\mathcal J_{6\to8}\). This
equality fixes the normalization dictionary that inserts
\(\gamma_{\rm HMT}=\Gamma_{6\to8}\) into the
Ashtekar--Barbero connection.

The information retained by reduction to the region is published through the
modular Hamiltonian. The starting structure of this reading is declared to be
\[
 \Delta_4^{\rm mod}
 =\frac{(A-C^*)\pi/180}{9\cdot90}
 \left(1-\frac7{12}\frac\pi{729}\right),
 \qquad w_{90}=90\Delta_4^{\rm mod},\quad
 w_{120}=120\Delta_4^{\rm mod},
\]
we define
\[
 K_{\rm link}(w)=wN,
 \qquad
 \rho_{\rm link}(w)=
 \frac{\operatorname{diag}(1,e^{-w},e^{-2w})}
 {1+e^{-w}+e^{-2w}},
\]
\begin{equation}
 \rho_A=
 \bigotimes_{e\in\partial A_{90}}\rho_{\rm link}(w_{90})
 \otimes
 \bigotimes_{e\in\partial A_{120}}\rho_{\rm link}(w_{120}),
 \qquad K_A=-\log\rho_A.
 \label{p12-83-c14-s2:eq:estado-hamiltoniano-modular-barbero-rev11}
\end{equation}
Then
\[
 K_A=\sum_{e\in\partial A}w_eN_e+c\mathbf1,
 \qquad
 S(\rho_A)=18S(w_{90})+18S(w_{120}),
\]
and the evaluation of that starting structure gives
\[
 S(\rho_A)=39.54837320881808\ldots\ \text{nats}
 =57.05624190358778\ldots\ \text{bits}.
\]
The quantity \(S(\rho_A)\) is the von Neumann entropy of the declared mixed
product. A global purification identifies it with an entanglement entropy
across the cut; without that purification, its canonical type is
reduced modular entropy. Thus, \(\gamma_{\rm HMT}\) fixes the quantum of
area/torsion, \(S_{\rm tri}\) counts the areal capacity, \(S_{\rm bi}\)
quantifies the allocation between sheets, and \(S(\rho_A)\) measures modular mixing.
Area operator, spectrum, and intertwiner belong to a single boundary
construction.
